We now proceed with the proof of Theorem~\ref{thm:localizedcase}. We begin with a series of three lemmas. Inspired by the results of \cite{dMGPSS20} in the case $E = \emptyset$, we consider the \emph{internal edge graph} of a subcomplex $\Delta \subset \lk_{\Gamma}(E)$. This is the graph contained in the $1$-skeleton of $\lk_{\Gamma}(E)$ consisting of edges $e \subset \Delta$ with $e \sqcup E$ interior. 

\begin{lemma}\label{lem:intedge}
Assume $\sigma(E)$ has codimension at least $2$. Let $\Delta$ be a subcomplex of $\lk_{\Gamma}(E)$, and assume $\Delta$ has no vertices $v$ with $\{v\} \sqcup E$ interior. If $L(\Gamma, E)|_{\Delta}$ is zero in degree $2$, then each connected component of the internal edge graph of $\Delta$ satisfies one of the following.
\begin{enumerate}
\item The component is a tree, and it has at most one  vertex $v$ with $\{v\} \sqcup E$ having carrier codimension more than $1$. 
\item The component has a unique cycle, and the carrier codimension of $\{w\} \sqcup E$ is equal to $1$ for every vertex $w$ in the component. 
\end{enumerate}
\end{lemma}

\begin{proof}
From \eqref{e:restrict}, we have the following exact sequence for the degree $2$ part of $L(\Gamma, E)|_{\Delta}$.
$$ \bigoplus_{\vert S \vert = 1} (I_S)_1\otimes_{k[\lk_{\Gamma}(E)]} k[\Delta]  \to I_2\otimes_{k[\lk_{\Gamma}(E)]} k[\Delta] \to (L(\Gamma, E)|_{\Delta})_2 \to 0.$$
Because $(L(\Gamma, E)|_{\Delta})_2 = 0$, the first map in the above complex is surjective.  As $\Delta$ has no vertices $v$ with $\{v\} \sqcup E$ interior, we see that
\begin{equation}\label{eq:intedge}
(x^{e}: e \subset \Delta, \enskip e \sqcup E \text{ is interior })_2 = (x^{\{v\}} \theta_{i}: v \subset \Delta, \enskip \sigma(\{v\} \sqcup E) = \{v_i\}^c)_2.
\end{equation}
Thus the number of edges $e$ with $e \sqcup E$ interior is less than or equal to the number of vertices $w$ with the carrier codimension of $\{w\} \sqcup E$ equal to $1$. If $\sigma(\{v\} \sqcup E) = \{v_i\}^c$ and $\theta_{i} = \sum_{w_j} a_{i,j}x^{\{w_j\}}$, then 
$$x^{\{v\}} \theta_{i} = \sum_{\{v, w_j\} \sqcup E \text{ interior }} a_{i,j}x^{\{v, w_j\}}.$$ In particular, both vector spaces in (\ref{eq:intedge}) naturally decompose into a direct sum of vector spaces indexed by the connected components of the internal edge graph. Therefore, in each connected component of the internal edge graph, the number of edges $e$ with $e \sqcup E$ interior is less than or equal to the number of vertices $v$ with $\{v\} \sqcup E$ of carrier codimension $1$. As the only connected graphs $(V, E)$ where $\vert E \vert \le \vert V \vert$ are either trees or contain a unique cycle, the result follows.
\end{proof}


\begin{lemma}\label{lem:nofourcycle}
Assume $\sigma(E)$ has codimension at least $2$. Let $F \subset \lk_{\Gamma}(E)$ be a face. Assume $F$ has no vertices $v$ with $\{v\} \sqcup E$ interior. If $L(\Gamma, E)|_{F}$ is zero in degree $2$, then no component of the internal edge graph of $F$ contains a cycle of length $4$. 
\end{lemma}
\begin{proof}
Suppose a component of the internal edge graph contains a $4$-cycle of vertices $F = \{t_1, t_2, u_1, u_2\}$. By Lemma~\ref{lem:intedge}, this is the unique cycle in this component and every vertex $w \in F$ has $\{w\} \sqcup E$ of carrier codimension $1$. Because $F$ is a face and there are no $3$-cycles in this component of the internal edge graph, we may assume that $\sigma(\{t_i\} \sqcup E) = \{v_1\}^c$ and $\sigma(\{u_i\} \sqcup E) = \{v_2\}^c$. Restricting to $F$ and using that $(L(\Gamma, E)|_{F})_2 = 0$, we have that 
$$(x^{\{ t_1, u_1 \}}, x^{ \{ u_1, t_2\}}, x^{\{t_2,u_2\}}, x^{\{u_2, t_1\}}) = (x^{\{t_1\}} \theta_{2}, x^{\{t_2\}} \theta_{2}, x^{\{u_1\}} \theta_{1}, x^{\{u_2\}} \theta_{1}).$$
The relation $\theta_{1} \theta_{2} - \theta_{2} \theta_{1} = 0$ expands into a relation between the generators of the right-hand side. But the left-hand side is $4$-dimensional, a contradiction.
\end{proof}

\begin{lemma}\label{lem:cd1case}
Assume $\sigma(E)$ has codimension $1$. Let $\Delta \subset \lk_{\Gamma}(E)$ be a subcomplex. Then 
$$\dim (L(\Gamma, E)|_{\Delta})_1 \ge |\{v \in \Delta: \{v\} \sqcup E \text{ interior}\}| - 1.$$
\end{lemma}
\begin{proof}
By considering the degree $1$ part of \eqref{e:restrict}, as the codimension of $\sigma(E)$ is $1$, 
we get the following exact sequence.
\begin{center}
\begin{tikzcd}[column sep = large]
 k \arrow[r] & \bigoplus\limits_{\mathclap{\substack{w \in \Delta \\ \{w\} \sqcup E \text{ interior }}}} k \cdot x^w \arrow[r] &   (L(\Gamma, E)|_{\Delta})_1  \arrow[r] & 0,
\end{tikzcd}
\end{center}
and the result follows. 
\end{proof}
